\clearpage
\begin{summarybox}
\begin{center}\begin{tikzpicture}[font=\small,>=stealth,line join=round]
\draw[gray!40,rounded corners=3pt] (0,0) rectangle (14.2,8.3);\node at (7.1,7.700000000000001) {三种写法，对应同一条抛物线};
\begin{scope}[shift={(1.7000000000000002,2.9642857142857144)},x=1.2cm,y=0.6571428571428571cm]\draw[->,gray] (-0.5,0)--(4.5,0) node[right]{$x$};\draw[->,gray] (0,-2)--(0,5) node[above]{$y$};\draw[blue,thick,solid,domain=-0.2:4.2,samples=180] plot(\x,{(\x-2)^2-1});\fill[orange!85!black] (0,3) circle(2.2pt);\fill[orange!85!black] (1,0) circle(2.2pt);\fill[orange!85!black] (3,0) circle(2.2pt);\fill[orange!85!black] (2,-1) circle(2.2pt);\draw[violet,dashed] (2,-.95)--(2,4.5);\node[left] at (0,3){$(0,3)$};\node[below=5pt] at (1,0){$1$};\node[below=5pt] at (3,0){$3$};\node[below=4pt] at (2,-1){$(2,-1)$};\end{scope}\node[anchor=west,align=left,text width=5.0cm] at (8.4,4.0) {$x^2-4x+3$\\[9pt]$=(x-2)^2-1$\\[9pt]$=(x-1)(x-3)$\\[9pt]顶点 $(2,-1)$\\[9pt]交点 $(1,0),(3,0)$};
\node[text width=12.6cm,align=center] at (7.1,.65) {一般式读纵轴交点，顶点式读顶点，两根式读横轴交点。};
\end{tikzpicture}\end{center}


\textbf{\textcolor{CSummary}{一句话带走}}
\textbf{给系数，用一般式；给顶点，用顶点式；给与 $x$ 轴的交点，用两根式。}
三种式子可能描述的是\textbf{同一条}抛物线，只是把不同的信息放在最显眼的位置。

\medskip
\textbf{\textcolor{CSummary}{三张最小记忆卡}}
\begin{itemize}[leftmargin=*,itemsep=0.4em,topsep=0.3em]
\item \textbf{一般式：}$y=ax^2+bx+c$，$a\ne0$；经过 $(0,c)$。
\item \textbf{顶点式：}$y=a(x-h)^2+k$；顶点 $(h,k)$、对称轴 $x=h$，适用于所有二次函数。
\item \textbf{两根式：}$y=a(x-x_1)(x-x_2)$；只有\textbf{存在实数根}（$\Delta\ge0$）时才能这样写实数的 $x_1,x_2$。
\end{itemize}

\medskip
\textbf{\textcolor{CSummary}{三条很值得记住的连线}}
\[
\boxed{h=-\frac{b}{2a}=\frac{x_1+x_2}{2}},\qquad
\boxed{k=\frac{4ac-b^2}{4a}},\qquad
\boxed{\Delta=b^2-4ac=-4ak}.
\]
第一条中“两个根的中点”部分需存在实数根；其余公式对所有实系数二次函数成立。

\medskip
\textbf{\textcolor{CSummary}{脑中留住这一幅图}}
\[
\underbrace{x^2-4x+3}_{\text{看 }(0,3)}
=\underbrace{(x-2)^2-1}_{\text{看顶点 }(2,-1)}
=\underbrace{(x-1)(x-3)}_{\text{看两个交点}}.
\]
两个交点 $1$、$3$ 的中点是 $2$，刚好是对称轴；最低点在 $x$ 轴下方，因此有两个交点。

\medskip
\textbf{\textcolor{CSummary}{考场30秒检查}}
先看 $a\ne0$；要判断交点个数看 $\Delta$ 或看顶点与 $x$ 轴的位置；要找顶点先配方；若已知两个交点和一个额外条件，就先设两根式只求 $a$。
\end{summarybox}
